\documentclass[11pt]{article}
\usepackage[margin=1in]{geometry}
\usepackage{amsmath,amssymb,amsthm}
\usepackage{hyperref}

\theoremstyle{plain}
\newtheorem{theorem}{Theorem}
\newtheorem{proposition}[theorem]{Proposition}
\newtheorem{corollary}[theorem]{Corollary}
\newtheorem{lemma}[theorem]{Lemma}
\theoremstyle{definition}
\newtheorem{definition}[theorem]{Definition}
\theoremstyle{remark}
\newtheorem{remark}[theorem]{Remark}

\newcommand{\dd}{\mathrm{d}}

\title{The Schwarzschild--de~Sitter family as one cubic on the de~Sitter throat:\\
\large a triple-angle chart, a dimension-selection theorem, and the family's description groupoid}
\author{Daryl Janzen\\
\small Department of Physics and Engineering Physics, University of Saskatchewan,\\
\small 116 Science Place, Saskatoon, SK S7N 5E2, Canada\\
\small \texttt{daryl.janzen@usask.ca}}
\date{}

\begin{document}
\maketitle

\begin{abstract}
The Schwarzschild--de~Sitter (Kottler) family with fixed cosmological constant is usually read as a
one-parameter family of distinct spacetimes, with de~Sitter and Nariai at its two ends. We study it instead as
one object: the horizon cubic of the family, with the de~Sitter curvature radius $\alpha=\sqrt{3/\Lambda}$ held
fixed and each member designated by one of its own horizon radii $r_0$, so that $2M/\alpha=x_0-x_0^{3}$ with
$x_0=r_0/\alpha$. Four groups of results follow. (i) The horizon locus of the whole family in the
$(x_0,x)$ plane is a line and an ellipse tilted at $45^{\circ}$, meeting exactly at the two Nariai
configurations; the exchange of which root is designated is an explicit involution of the line element, with
Nariai its unique fixed point. (ii) Writing $x_0=\varrho\sin w$, the horizon relation becomes a single
harmonic, $2M/\alpha=\tfrac{2}{3\sqrt3}\sin 3w$, for exactly one scale, $\varrho=2/\sqrt3$; this is Vi\`ete's
solution of the cubic, here with its angle identified as a sky angle and its scale as the gnomonic image of the
de~Sitter throat. In $D$ spacetime dimensions the analogous single-harmonic form exists if and only if $D=4$ or
$D=5$, and has zero mean only for $D=4$; the proof is a closed-form coefficient condition,
$\varrho^{2}=4j(D-1-j)/((D-1)(D-2))$ for the harmonic $(D-1-2j)w$. (iii) On the equatorial section of the
de~Sitter hyperboloid, the planar figure that carries this chart---three hinges at distance $2\alpha$ about a
throat of radius $\alpha$---is fixed by the throat alone, and the sections of the hyperboloid by the slicing
planes realise the horizon trichotomy as a conic trichotomy, the Nariai section being a pair of the
hyperboloid's null rulings. (iv) The family's descriptions form a groupoid generated by the root exchange
$\sigma$ and the sky-angle step $\tau$, subject to $(\sigma\tau)^{2}=1$; its monodromy over the mass plane is
$S_{3}$, the Galois group of the horizon cubic, and its full discrete symmetry is $\mathrm{Aut}(A_{2})=S_{3}\times
\mathbb{Z}_{2}$, the extra $\mathbb{Z}_{2}$ being the mass reflection. We state precisely which of these are
statements about an embedding and which are identities of the cubic read on the figure; no member with $M\neq0$
is claimed to embed in de~Sitter space.
\end{abstract}

\section{Introduction}\label{sec:intro}

With $\Lambda>0$ the static spherically symmetric vacuum solutions of Einstein's equations are the
Schwarzschild--de~Sitter, or Kottler, metrics~\cite{Kottler1918},
\begin{equation}\label{eq:sds}
\dd s^{2}=-f(r)\,\dd t^{2}+\frac{\dd r^{2}}{f(r)}+r^{2}\dd\Omega^{2},\qquad
f(r)=1-\frac{2M}{r}-\frac{r^{2}}{\alpha^{2}},\qquad \alpha^{2}=\frac{3}{\Lambda}.
\end{equation}
For $0<27M^{2}<\alpha^{2}$ the metric has a black-hole horizon and a cosmological horizon; at $27M^{2}=\alpha^{2}$
they merge, and the near-horizon region of the degenerate member is the Nariai geometry
$\mathrm{dS}_{2}\times S^{2}$~\cite{Nariai1951,GinspargPerry1983}; for larger $M$ there is no static region. The
horizon radii are the positive roots of the cubic $rf(r)=0$, and their trigonometric form is in the literature:
Visser writes the event horizon as $r_{E}=2a\sin\bigl(\tfrac13\arcsin(3M/a)\bigr)$ with $a=\alpha/\sqrt3$, and
carries the cubic's negative root as a formal ``virtual'' horizon in area-product relations~\cite{Visser2013}.
The static radius $r=(M\alpha^{2})^{1/3}$, where the radial acceleration of a static observer vanishes, and the
photon sphere $r=3M$ are classical~\cite{StuchlikHledik1999}.

This paper reads the same family as a single object. Two choices make that possible. First, $\alpha$ is held
fixed and the family is parametrised not by $M$ but by one of the family's own horizon radii, so that
designating a root is the basic operation and the mass is a derived quantity. Second, the parameter is written
as the sine of an angle, and the scale of that sine is chosen so that the horizon relation is a single harmonic.
Both choices are elementary. What they return is not: a planar locus for the whole family whose singular points
are the Nariai configurations; an explicit involution of the line element exchanging two designations of one
geometry; a scale that the single-harmonic condition fixes uniquely, and which is the radius of the gnomonic
image of the de~Sitter throat seen from a definite distance; a theorem that the single-harmonic form exists in
four and five spacetime dimensions and in no other; a figure on the equatorial section of de~Sitter space that
carries the chart and realises the horizon trichotomy as a trichotomy of conic sections, with the degenerate
member cut out by two null rulings; and the family's discrete symmetry as a groupoid whose monodromy is the
Galois group of the horizon cubic.

We are careful about one thing throughout. The figure lives on the de~Sitter hyperboloid, but the members of the
family with $M\neq0$ are not sections of it: every hyperplane section of the hyperboloid has constant curvature
(a section by a hyperplane with spacelike normal is a de~Sitter space of smaller radius), and the Kretschmann scalar of~\eqref{eq:sds}, $48M^{2}/r^{6}+24/\alpha^{4}$, is constant only at
$M=0$. Every statement below is therefore either an identity of the horizon cubic, a statement about a defined
planar figure, or a statement about sections of the hyperboloid, and we say which.

Section~\ref{sec:defs} sets out the objects. Section~\ref{sec:locus} treats the horizon locus and the root
exchange. Section~\ref{sec:triple} proves the single-harmonic theorem and its dimension selection.
Sections~\ref{sec:figure}--\ref{sec:groupoid} treat the figure, the slicing surface and the groupoid.

\section{The objects}\label{sec:defs}

Throughout, lengths are measured in units of $\alpha$: $x=r/\alpha$, and $\mu=2M/\alpha$.

\begin{definition}[The family and its horizon cubic]\label{def:family}
For fixed $\alpha>0$ the \emph{family} is the set of metrics~\eqref{eq:sds} with $M\in\mathbb{R}$. The
\emph{horizon cubic} of the member $\mu$ is
\begin{equation}\label{eq:cubic}
p_{\mu}(x)=x^{3}-x+\mu,
\end{equation}
whose real roots are the zeros of $xf$; its positive roots are the member's horizons.
\end{definition}

The cubic has vanishing $x^{2}$ coefficient, so its three roots, real or complex, sum to zero. Its discriminant
is $\Delta(\mu)=4-27\mu^{2}$: three distinct real roots for $|\mu|<2/(3\sqrt3)$ (the \emph{undercritical}
members), a double root at $|\mu|=2/(3\sqrt3)$ (the \emph{Nariai} members, $\Lambda M^{2}=1/9$), and one real
root for larger $|\mu|$ (the \emph{overcritical} members). No member has a triple root: three equal roots $a$
would require $3a^{2}=1$ and $a^{3}=0$ at once.

\begin{definition}[Designation]\label{def:designation}
A \emph{designation} of a member is a choice of one of its real horizon-cubic roots, $x_{0}$. Since $x_{0}$ is a
root,
\begin{equation}\label{eq:designation}
\mu=x_{0}-x_{0}^{3}.
\end{equation}
\end{definition}

Equation~\eqref{eq:designation} is a parametrisation of the family by a root, and it is many-to-one: an
undercritical member is met at three values of $x_{0}$, one per root. The family is thereby parametrised with
the mass a derived quantity and $\alpha$ the only fixed scale, and the three-fold covering is the first
structure the rest of the paper reads.

\begin{definition}[The horizon locus]\label{def:locus}
The \emph{horizon locus} is the set of pairs $(x_{0},x)\in\mathbb{R}^{2}$ such that $x$ is a horizon-cubic root of
the member designated by $x_{0}$; equivalently, $x^{3}-x=x_{0}^{3}-x_{0}$.
\end{definition}

\begin{definition}[The sky-angle chart]\label{def:chart}
For a scale $\varrho>0$, the \emph{sky-angle chart} of scale $\varrho$ is the parametrisation
$x_{0}=\varrho\sin w$, $w\in\mathbb{R}$.
\end{definition}

The name is explained in Section~\ref{sec:figure}, where $w$ is the angle between a line of sight and the centre
of the throat's image on a celestial sphere and $\varrho$ the radius of that image in the gnomonic projection;
nothing in Sections~\ref{sec:locus} and~\ref{sec:triple} uses that reading.

\section{The horizon locus and the root exchange}\label{sec:locus}

\begin{proposition}[Factorisation and the locus]\label{prop:locus}
With $\mu=x_{0}-x_{0}^{3}$,
\begin{equation}\label{eq:factor}
p_{\mu}(x)=(x-x_{0})\bigl(x^{2}+x\,x_{0}+x_{0}^{2}-1\bigr),
\end{equation}
so the horizon locus is the union of the line $x=x_{0}$ and the conic
\begin{equation}\label{eq:ellipse}
E:\quad x^{2}+x\,x_{0}+x_{0}^{2}=1 .
\end{equation}
$E$ is an ellipse with semi-axes $\sqrt2$ along the anti-diagonal $x=-x_{0}$ and $\sqrt{2/3}$ along the diagonal
$x=x_{0}$. The line and the ellipse meet exactly at $(x_{0},x)=\pm(1/\sqrt3,1/\sqrt3)$, the Nariai
configurations with the merged root designated, and the vertical tangents of $E$ are at $x_{0}=\pm2/\sqrt3$, the
Nariai configurations with the lone root designated.
\end{proposition}

\begin{proof}
$p_{\mu}(x)=(x^{3}-x_{0}^{3})-(x-x_{0})$, and $x^{3}-x_{0}^{3}=(x-x_{0})(x^{2}+xx_{0}+x_{0}^{2})$. The quadratic form
$x^{2}+xx_{0}+x_{0}^{2}$ has matrix $\bigl(\begin{smallmatrix}1&1/2\\1/2&1\end{smallmatrix}\bigr)$, eigenvalues
$1/2$ and $3/2$ with eigenvectors $(1,-1)$ and $(1,1)$, so $E$ is an ellipse with the stated semi-axes.
Substituting $x=x_{0}$ in~\eqref{eq:ellipse} gives $3x_{0}^{2}=1$. Solving~\eqref{eq:ellipse} for $x$ gives
$x=\bigl(-x_{0}\pm\sqrt{4-3x_{0}^{2}}\bigr)/2$, real exactly for $|x_{0}|\le2/\sqrt3$, with the two branches meeting
at $|x_{0}|=2/\sqrt3$, where $\mu=x_{0}-x_{0}^{3}=\mp2/(3\sqrt3)$.
\end{proof}

The ellipse is not new as a curve: it appears, with the factorisation, in the author's dissertation~\cite{JanzenThesis}.
What the proposition adds is its place in the family: the degenerate members are the singular points of the
family's own locus, where its two components cross, and the end points of its minor axis. The ellipse has a
second reading that will be used in Section~\ref{sec:groupoid}: since the roots sum to zero, the triple
$(x_{1},x_{2},x_{3})$ lies in the plane $x_{1}+x_{2}+x_{3}=0$, and parametrising that plane by $(x_{1},x_{2})$ the
induced Euclidean norm is $2(x_{1}^{2}+x_{1}x_{2}+x_{2}^{2})$. So $E$ is the circle of radius $\sqrt2$ of the
Euclidean metric on the sum-zero plane---the plane in which the root system $A_{2}$ is usually drawn.

\begin{proposition}[The root exchange]\label{prop:sigma}
Let
\begin{equation}\label{eq:sigma}
\sigma(x_{0})=\tfrac12\Bigl(-x_{0}+\sqrt{4-3x_{0}^{2}}\Bigr).
\end{equation}
Then $\sigma(x_{0})$ is a root of the member designated by $x_{0}$, so $\sigma(x_{0})-\sigma(x_{0})^{3}=x_{0}-x_{0}^{3}$
and the metric function $f$ is unchanged when $\sigma(x_{0})$ is designated in place of $x_{0}$. On
$[-1/\sqrt3,2/\sqrt3]$, $\sigma$ is a decreasing involution exchanging $[0,2/\sqrt3]$ with $[-1/\sqrt3,1]$;
it maps $[0,1]$ onto itself, exchanging $0$ and $1$, and its unique fixed point is $x_{0}=1/\sqrt3$.
\end{proposition}

\begin{proof}
$\sigma(x_{0})$ is the upper branch of $E$ over $x_{0}$, hence a root by Proposition~\ref{prop:locus}, and $f$
depends on the designation only through $\mu$. Put $s=\sqrt{4-3x_{0}^{2}}$. Then
$4-3\sigma(x_{0})^{2}=\bigl((3x_{0}+s)/2\bigr)^{2}$, and $3x_{0}+s\ge0$ exactly on $[-1/\sqrt3,2/\sqrt3]$, so there
$\sqrt{4-3\sigma(x_{0})^{2}}=(3x_{0}+s)/2$ and $\sigma(\sigma(x_{0}))=\tfrac12\bigl(-\tfrac12(-x_{0}+s)+\tfrac12(3x_{0}+s)\bigr)
=x_{0}$. A fixed point satisfies $s=3x_{0}$, so $x_{0}=1/\sqrt3$. The derivative is $-(3x_{0}+s)/(2s)<0$ on the
interior, and $\sigma(0)=1$, $\sigma(1)=0$, $\sigma(2/\sqrt3)=-1/\sqrt3$.
\end{proof}

So two designations of one geometry are exchanged by an explicit map that leaves the line element unchanged.
At $x_{0}=0$ the member is de~Sitter ($\mu=0$, roots $0,\pm1$), and $\sigma$ carries it to the designation $x_{0}=1$
of the same geometry; Nariai is where the exchange has nothing to exchange. The third root, the negative one for
$0<\mu<2/(3\sqrt3)$, is reached from the other two not by $\sigma$ but by the reflection
\begin{equation}\label{eq:R}
R:(x_{0},x)\mapsto(-x_{0},-x),\qquad \mu\mapsto-\mu,
\end{equation}
which preserves both components of the locus and changes the sign of the mass. $R$ is not a symmetry of a
member; it is a symmetry of the family, carrying the member $\mu$ to the member $-\mu$.

\section{The single-harmonic chart and the dimension selection}\label{sec:triple}

\begin{theorem}[The triple angle]\label{thm:triple}
In the sky-angle chart of scale $\varrho$,
\[
x_{0}-x_{0}^{3}=\bigl(\varrho-\tfrac34\varrho^{3}\bigr)\sin w+\tfrac14\varrho^{3}\sin 3w ,
\]
so the horizon relation is a single harmonic of $w$ if and only if $\varrho=2/\sqrt3$, in which case
\begin{equation}\label{eq:triple}
\mu=\frac{2}{3\sqrt3}\,\sin 3w .
\end{equation}
The three roots of the member $\mu$ are then the designations at $w$, $\pi/3-w$ and $-\pi/3-w$, and the
reflection $w\mapsto\pi/3-w$ is the root exchange $\sigma$ of Proposition~\ref{prop:sigma}, with fixed point the
crest $w=\pi/6$ of $\sin3w$.
\end{theorem}

\begin{proof}
The expansion is $\sin^{3}w=\tfrac14(3\sin w-\sin3w)$. The $\sin w$ coefficient vanishes for positive $\varrho$
only at $\varrho^{2}=4/3$, and then $\tfrac14\varrho^{3}=2/(3\sqrt3)$. Since $\sin3w$ takes the same value at $w$,
$\pi/3-w$ and $-\pi/3-w$, those three designations carry the same $\mu$, and they are distinct roots for
$w\notin\tfrac{\pi}{6}+\tfrac{\pi}{3}\mathbb{Z}$. With $x_{0}=\tfrac{2}{\sqrt3}\sin w$ one has
$\tfrac{2}{\sqrt3}\sin(\tfrac{\pi}{3}-w)=\cos w-\tfrac{1}{\sqrt3}\sin w$, and
$\sigma(x_{0})=\tfrac12\bigl(-\tfrac{2}{\sqrt3}\sin w+2\cos w\bigr)$ for $w\in[-\pi/6,\pi/2]$, which is the same.
\end{proof}

Equation~\eqref{eq:triple} solved for $x_{0}$ is Vi\`ete's trigonometric solution of the depressed cubic, and in
that form it is in the literature~\cite{Visser2013}. What the theorem adds is that the scale is forced by the
single-harmonic condition, and the identification, in Section~\ref{sec:figure}, of $w$ as a sky angle and
$\varrho=2/\sqrt3$ as the radius of a gnomonic image of the throat. The single-harmonic condition also selects
a dimension.

\begin{theorem}[Dimension selection]\label{thm:dimension}
In $D\ge4$ spacetime dimensions the Schwarzschild--Tangherlini--de~Sitter metric function is
$f=1-2M/r^{D-3}-r^{2}/\alpha^{2}$, and the horizon relation at a designated root is
$\mu_{D}=x_{0}^{D-3}-x_{0}^{D-1}$. In the sky-angle chart of scale $\varrho$, the harmonic $(D-1-2j)w$, for
$1\le j<(D-1)/2$, has vanishing coefficient if and only if
\begin{equation}\label{eq:rhoj}
\varrho^{2}=\frac{4j\,(D-1-j)}{(D-1)(D-2)} .
\end{equation}
Consequently a scale exists for which $\mu_{D}$ is a single harmonic of $w$ plus a constant if and only if $D=4$
or $D=5$. The scale is then unique, $\varrho^{2}=4/(D-1)$, and the result has zero mean only for $D=4$:
\[
D=4:\ \mu_{4}=\frac{2}{3\sqrt3}\sin3w;\qquad D=5:\ \mu_{5}=\frac18\bigl(1-\cos4w\bigr).
\]
\end{theorem}

\begin{proof}
Write $n=D-1$. For $m\ge1$,
$\sin^{m}w=2^{1-m}\sum_{k=0}^{\lfloor (m-1)/2\rfloor}(-1)^{\lfloor m/2\rfloor-k}\binom{m}{k}\,h_{m-2k}(w)$ plus, for even
$m$, the constant $2^{-m}\binom{m}{m/2}$, where $h_{\ell}=\sin\ell w$ for odd $m$ and $\cos\ell w$ for even $m$.
The harmonic $h_{n-2j}$ appears in $-\varrho^{n}\sin^{n}w$ with $k=j$ and in $\varrho^{n-2}\sin^{n-2}w$ with $k=j-1$,
and the two sign factors $(-1)^{\lfloor n/2\rfloor-j}$ and $(-1)^{\lfloor (n-2)/2\rfloor-(j-1)}$ are equal. The
coefficient therefore vanishes if and only if $\varrho^{n-2}2^{3-n}\binom{n-2}{j-1}=\varrho^{n}2^{1-n}\binom{n}{j}$,
which is~\eqref{eq:rhoj} since $\binom{n-2}{j-1}/\binom{n}{j}=j(n-j)/(n(n-1))$. The function $j(n-j)$ increases
strictly for $j<n/2$, so the conditions for two distinct non-constant lower harmonics are incompatible; one
scale clears all of them only when there is exactly one, $j=1$ with $n-2\ge1>n-4$, that is $n\in\{3,4\}$.
Then $\varrho^{2}=4/n$. For $n=3$ there is no constant term; for $n=4$ the constant survives, and direct
substitution gives the two closed forms.
\end{proof}

\begin{remark}
The metric function used in the theorem is the standard Tangherlini--de~Sitter form. It is also the entire
solution space of the first-order linear equation $rf'+(D-3)(f-1)+2\Lambda r^{2}/(D-2)=0$ which the vanishing of
$G^{t}{}_{t}+\Lambda$ imposes on a metric of the form~\eqref{eq:sds} with $S^{D-2}$ in place of $S^{2}$, with $M$
the constant of integration and $\alpha^{2}=(D-1)(D-2)/2\Lambda$.
\end{remark}

The theorem is a statement about the family and a chart, and it is sharp: in six or more dimensions the
designations $2\pi/(D-1)$ apart in $w$ no longer carry a common mass for any scale, and the three-fold (in four
dimensions) or four-fold (in five) structure of the dial has no analogue. The difference between the two
surviving dimensions is the parity of $\mu_{D}$ in $x_{0}$: it is odd for even $D$ and even for odd $D$, so the
reflection $R$ of~\eqref{eq:R} relates distinct members only in four dimensions, while in five it fixes each.

\section{The figure on the de~Sitter throat}\label{sec:figure}
\emph{[To be drafted: the equatorial section $-X_{0}^{2}+X_{1}^{2}+X_{2}^{2}=\alpha^{2}$; the throat circle; the
gnomonic projection and the charting distance $\sqrt7\,\alpha/2$ at which the throat's image has radius
$2/\sqrt3$; the hinge at $2\alpha$ fixed by the throat (subtense $60^{\circ}$, incircle, nine-point circle,
tangent length equal to height); the plane sections of the hyperboloid at transverse distance $d$---hyperbola,
two null rulings, secant hyperbola---matching the horizon trichotomy; the Nariai section as the two rulings; and
the statement that sections are de~Sitter, so that no member with $M\neq0$ is a section.]}

\section{The slicing surface}\label{sec:surface}
\emph{[To be drafted: the surface $\epsilon\,\dd l^{2}+r(l)^{2}\dd\phi^{2}$ with $\dd r/\dd l=\sqrt{|f|}$; Gaussian
curvature $1/\alpha^{2}-M/r^{3}$ on every piece, blind to the horizons; its sign change at the static radius,
which is also the zero of the comoving acceleration; the continuation $\theta\mapsto\pi/2+i\psi$ through the
equatorial seam with the signature flip confined to the surface; the curve through $r=0$ as a branch point onto
the conjugate root.]}

\section{The description groupoid}\label{sec:groupoid}
\emph{[To be drafted: the groupoid of designations with generators $\sigma$ and $\tau:w\mapsto w+2\pi/3$, the
relation $(\sigma\tau)^{2}=1$ and its completeness; the branched cover of the $\mu$-line with monodromy $S_{3}$
equal to the Galois group of $p_{\mu}$ over $\mathbb{Q}(\mu)$; $\mathrm{Aut}(A_{2})=S_{3}\times\mathbb{Z}_{2}$ with $R$
the outer factor; the three distinct operations $\sigma$, $R$ and the seam continuation; Weinstein's groupoid
framework~\cite{Weinstein1996}.]}

\section{Discussion}\label{sec:discussion}
\emph{[To be drafted.]}

\begin{thebibliography}{99}
\bibitem{Kottler1918} F.~Kottler, \emph{\"Uber die physikalischen Grundlagen der Einsteinschen
Gravitationstheorie}, Ann.\ Phys.\ (Leipzig) \textbf{361} (1918) 401.
\bibitem{Nariai1951} H.~Nariai, \emph{On a new cosmological solution of Einstein's field equations of
gravitation}, Sci.\ Rep.\ T\^ohoku Univ.\ Ser.\ I \textbf{35} (1951) 62; reprinted in Gen.\ Rel.\ Grav.\
\textbf{31} (1999) 963.
\bibitem{GinspargPerry1983} P.~Ginsparg and M.~J.~Perry, \emph{Semiclassical perdurance of de~Sitter space},
Nucl.\ Phys.\ B \textbf{222} (1983) 245.
\bibitem{Visser2013} M.~Visser, \emph{Area products for stationary black hole horizons}, Phys.\ Rev.\ D
\textbf{88} (2013) 044014.
\bibitem{StuchlikHledik1999} Z.~Stuchl\'ik and S.~Hled\'ik, \emph{Some properties of the
Schwarzschild--de~Sitter and Schwarzschild--anti-de~Sitter spacetimes}, Phys.\ Rev.\ D \textbf{60} (1999)
044006.
\bibitem{JanzenThesis} D.~Janzen, \emph{A Solution to the Cosmological Problem of Relativity Theory}, PhD
dissertation, University of Saskatchewan (2012), ch.~4.
\bibitem{Weinstein1996} A.~Weinstein, \emph{Groupoids: unifying internal and external symmetry}, Notices
Amer.\ Math.\ Soc.\ \textbf{43} (1996) 744.
\end{thebibliography}

\end{document}
